\chapter{Materials and Methods}

\section{Materials}

\subsection{Reagents}


\begin{longtable}{L{0.45\linewidth} R{0.45\linewidth}}
\rowcolor{white}
\caption{Reagents used in this Project}
\label{tab:reagents}
\\
\hline
\textbf{Reagent}& \textbf{Manufacturer}\\                           
\endfirsthead
\rowcolor{white}
\caption{Continuation}
\\
\hline
\textbf{Reagent}& \textbf{Manufacturer}\\ 
\endhead
\hline
\endfoot
Acetic acid 100\%       & AppliChem, Darmstadt, Germany          \\
Acrylamid & ??? \\
Agar                    & AppliChem,   Darmstadt, Germany        \\
Agarose LE              & Genaxxon, Ulm, Germany                 \\
Ammonium persulfate     & Merck,   Darmstadt, Germany            \\
Ampuwa                  & Fresenius, Bad Homburg v.d.H., Germany \\
APS & ??? \\
BSA          &                                                   \\
Bromphenol   blue       & Merck, Darmstadt, Germany              \\
Chloroform              & AppliChem,   Darmstadt, Germany        \\
EDTA disodium dihydrate & AppliChem,   Darmstadt, Germany        \\
Ethanol 100\%           & AppliChem,   Darmstadt, Germany        \\
Ethidium bromide 0.07\% & AppliChem,   Darmstadt, Germany        \\
Liquid   nitrogen       & Air Products, Allentown, USA           \\
Formaldehyde 37\%       & AppliChem,   Darmstadt, Germany        \\
Glucose                 & AppliChem,   Darmstadt, Germany        \\
Glycerin 87\%           & AppliChem, Darmstadt,   Germany        \\
Glycerin, H2O free      & AppliChem,   Darmstadt, Germany        \\
HEPES                   & AppliChem,   Darmstadt, Germany        \\
Isoamyl alcohol         & Riedel-de   Haen, Seelze, Germany      \\
Isopropyl alcohol       & AppliChem,   Darmstadt, German         \\
L-arabinose             & Sigma-Aldrich, St. Louis, USA          \\
LB Agar                 & Invitrogen,   Waltham, USA             \\
Lysozyme                & Sigma-Aldrich, St. Louis, US           \\
Methanol                & Fisher Scientific, Hampton, USA        \\
Paraformaldehyd         &                                         \\
Polyacrylamide          &                                        \\
Phenol                  & AppliChem,   Darmstadt, Germany        \\
Polybrene               & Sigma-Aldrich,   St. Louis, USA        \\
Potassium acetate       & AppliChem,   Darmstadt, Germany        \\
Potassium chlorate      & AppliChem,   Darmstadt, Germany        \\
SDS ultrapure           & AppliChem,   Darmstadt, Germany        \\
Silver nitrate          & AppliChem,   Darmstadt, Germany        \\
Sodium acetate          & AppliChem,   Darmstadt, Germany        \\
Sodium carbonate        & AppliChem,   Darmstadt, Germany        \\
Sodium chloride         & AppliChem,   Darmstadt, Germany        \\
Sodium hydroxide        & AppliChem,   Darmstadt, Germany        \\
Sodium thiosulfate      & AppliChem,   Darmstadt, Germany        \\
TEMED                   & AppliChem,   Darmstadt, Germany        \\
Tris ultrapure          & AppliChem,   Darmstadt, Germany        \\
Triton X-100            &                                        \\
Tween 20                & AppliChem,   Darmstadt, Germany        \\
Xylencyanol             & AppliChem,   Darmstadt, Germany      \\ 

\end{longtable}

\subsection{Enzymes and Buffers}

\begin{longtable}{L{0.45\linewidth} R{0.45\linewidth}}
\rowcolor{white}
\caption{Enzymes and Enzyme Buffers used in this Project}
\label{tab:Enzymes}
\\
\hline
\textbf{Product}& \textbf{Supplier}\\                           
\endfirsthead
\rowcolor{white}
\caption{Continuation}
\\
\hline
\textbf{Product}& \textbf{Supplier}\\ 
\endhead
\hline
\endfoot
10x buffer R & Thermo Scientific, Rockford, USA \\
10x r2.1 buffer & New England BioLabs, Ipswich, USA \\
BbsI restriction enzyme & New England BioLabs, Ipswich, USA \\
GoTaq buffer & Promega, Mannheim, Germany \\
GoTaq DNA Polymerase & Promega, Mannheim, Germany \\
PstI restriction enzyme & Thermo Scientific, Rockford, USA \\
Proteinase K & AppliChem, Darmstadt, Germany \\
Q5 buffer & New England BioLabs, Ipswich, USA \\
Q5 DNA Polymerase & New England BioLabs, Ipswich, USA \\
Ribonuklease A (RNAse) & AppliChem, Darmstadt, Germany \\
T4 ligase & \\
10x T4 buffer & \\
\end{longtable}

\subsection{Antibiotics}

\begin{longtable}{L{0.45\linewidth} R{0.45\linewidth}}
\rowcolor{white}
\caption{Antibiotics used in this Project}
\label{Antibiotics_table}
\\
\hline
\textbf{Antibiotic}& \textbf{Supplier}\\                           
\endfirsthead
\rowcolor{white}
\caption{Continuation}
\\
\hline
\textbf{---}& \textbf{---}\\ 
\endhead
\hline
\endfoot
Puromycin & \\
Streptomycin & Roth, Karlsruhe, Germany \\
\end{longtable}

\subsection{Appliances}

\begin{longtable}{L{0.45\linewidth} R{0.45\linewidth}}
\rowcolor{white}
\caption{Appliances used in this Project}
\label{tab:devices}
\\
\hline
\textbf{Appliance}& \textbf{Manufacturer}\\                           
\endfirsthead
\rowcolor{white}
\caption{Continuation}
\\
\hline
\textbf{Appliance}& \textbf{Manufacturer}\\ 
\endhead
\hline
\endfoot
Amersham nitrocellulose membrane 0.45\(\mu\)m & ?? \\
Avanti J-25 Centrifuge                    & Beckmann Coulter, Brea, USA          \\
Biological safety cabinet Hera Safe       & Heraeus, Hanau, Germany              \\
Biophotometer                             & Eppendorf, Hamburg, Germany          \\
Centrifuge 5415 R                         & Eppendorf, Hamburg, Germany          \\
Cytoflex FACS                             & Beckmann Coulter, Brea, USA          \\
Electrophoresis Power Supply              & ThermoScientific, Rockford, USA     \\
Eporator                                  & Eppendorf, Hamburg, Germany         \\
Gel Doc EZ Imager                         & BioRad, Hercules, USA                \\
Gel electrophoresis large gel box         & BioStep, Burkhardtsdorf Germany       \\
Gel electrophoresis small gel box         & BioStep, Burkhardtsdorf Germany       \\
Heraeus Incubator                         & Thermo Scientific, Rockford, USA      \\
Herafreeze Basic -80                      & Thermo Scientific, Rockford, USA      \\
Incubator Shaker G25                      & NBS, Edison, USA                      \\
Inolab pH-Meter                           & WTW, Weilheim, Germany                \\
Laboratory scale Sartorius                & Sartorius, Göttingen, Germany         \\
Leitz DMIL fluorescence microscope        & Leica, Wetzlar, Germany               \\
Leitz DMIL light microscope               & Leica, Wetzlar, Germany              \\
Magnet and magnetic stirrer RCT IKAMAG    & IKA, Staufen, Germany                 \\
Megafone 40 Centrifuge                    & Heraeus, Hanau, Germany               \\
Mikro 220R Centrifuge                     & Eppendorf, Hamburg Germany            \\
Milli-Q direct                            & Merck, Darmstadt, Germany            \\
Mixer HC thermomixer                      & StarLab, Brussels, Belgium           \\
improved Neubauer chamber & ??? \\
Odyssey CLx Imager & LI-COR Biosciences, Lincoln, USA\\
Optima XPN-80 Ultracentrifuge             & Beckman Coulter, Brea, USA           \\
PD-10 columns                             & GE Healthcare, Chicago, USA           \\
PowerPac 200 Electrophoresis Power Supply & BioRad, Hercules, USA                 \\
RC 5C Plus Centrifuge                     & ThermoScientific Rockford, USA       \\
Sartorius precision scale research        & Sartorius, Göttingen, Germany        \\
T Professional ThermoCycler (PCR)         & Biometra Analytik Jena, Germany      \\
Ultracentrifuge tubes 13.2 ml             & Beckman Coutler, Brea, USA           \\
Varioklav steam sterilizer                & H+P, Oberschleißheim, Germany        \\
VortexGenie 2                             & Scientific Industries, Bohemia, USA \\


\end{longtable}

\subsection{Dishes and Pipettes}
\begin{longtable}{L{0.45\linewidth} R{0.45\linewidth}}
\rowcolor{white}
\caption{Dishes and Pipettes used in this Project}
\label{Dishes and Pipettes used in this Project}
\\
\hline
\textbf{Appliance}& \textbf{Manufacturer}\\                           
\endfirsthead
\rowcolor{white}
\caption{Continuation}
\\
\hline
\textbf{Appliance}& \textbf{Manufacturer}\\ 
\endhead
\hline
\endfoot
Tissure culture dish, ØxH: 150mmx20mm &	Sarstedt, Nümbrecht, Germany \\
Tissure culture dish, ØxH: 150mmx20mm &	Sarstedt, Nümbrecht, Germany \\ 
Cell culture plate, 24-well	 & Sarstedt, Nümbrecht, Germany \\ 
Cell culture plate, 96-well	 & Sarstedt, Nümbrecht, Germany \\ 
Nunc Serological Pipette, 2, 5, 10, 25, 50ml & Sarstedt, Nümbrecht, Germany \\
\end{longtable}

\subsection{Cell lines}

\begin{longtable}{L{0.45\linewidth} R{0.45\linewidth}}
\rowcolor{white}
\caption{Cell lines used in this Project}
\label{tab:cellines}
\\
\hline
\textbf{Cell line}& \textbf{Manufacturer/ATCC number}\\                           
\endfirsthead
\rowcolor{white}
\caption{Continuation}
\\
\hline
\textbf{---}& \textbf{---}\\ 
\endhead
\hline
\endfoot

HEK-293 & ATCC CRL-1573TM, Virginia, USA \\
A549 & ATCC CCL-185, Virginia, USA \\

\end{longtable}

\subsection{Software}

\begin{longtable}{L{0.45\linewidth} R{0.45\linewidth}}
\rowcolor{white}
\caption{Software used in this Project}
\label{tab:software}
\\
\hline
\textbf{Software}& \textbf{Supplier}\\                           
\endfirsthead
\rowcolor{white}
\caption{Continuation}
\\
\hline
\textbf{Software}& \textbf{Supplier}\\ 
\endhead
\hline
\endfoot
ApE - a Plasmid Editor & Jorgensen Lab \\
Cytoflex & Beckmann Coulter, Brea, USA \\
ImageLab & BioRad, Hercules, USA \\
Overleaf, online \LaTeX \ editor & Overleaf c/o Digital Science, London, UK \\
\end{longtable}

\newpage

\section{Microbial Methods}

\subsection{Bacterial fluid cultures}

Only DH10\(\beta\) E. Coli cells were used in this project. Cells were stored at -80\C \ and thawed on ice in chilled \ce{ddH2O} before use. when growing in fluid culture, LB medium with the appropriate antibiotics was used. The culture temperature was 37 \C, shaking was done at 1,100 rpm if not stated otherwise.

\subsection{Bacterial culture on agar dishes}

To single-out a clonal bacterial population. Lb Agar plates with appropriate antibiotics were prepared and around 100 \ul of bacterial cell suspension, the exact amount varying slightly depending on observed colony density in previous attempts, was streaked on the plate with a sterile serological pipette. Plates were warmed before use and labeled to avoid mixups. the plates were incubated upside down and for a maximum of 16 hours.

\subsection{Small scale Plasmid DNA preparation}
For the small scale preparation (mini-prep) of plasmids, 1.8 ml DH10\(\beta\) overnight cultures were centrifuged at 16.500 x g and the bacterial pellet re-suspended by vortexing in 100 \ul GTE-Buffer. All of the following centrifuging steps were carried out at 16.500 x g. 200 \ul Lysis-Buffer were added and mixing was done by inverting the tube multiple times, the lysis reaction was neutralized after one minute by adding 150 \ul KAc solution and the tube was inverted a further 4-6 times. Centrifuging was done for 10 minutes at room temperature, the supernatant transferred to a new reaction tube and the DNA precipitated with 1 ml EtoH 100\% at -20\C. The centrifuging was repeated for 20 minutes at 4\C. The resulting DNA-pellet was washed with EtoH 70\% at -20\C \ and after a final centrifuging step for 10 minutes at 4\C \, the pellet was air-dryed for five minutes and taken up in 25 \ul TE-Buffer.

\begin{longtable}[ht]{L{0.45\linewidth} R{0.45\linewidth}}
\rowcolor{white}
\caption{Mini-prep Buffer composition}
\label{Mini-prep Buffer}
\\
\hline
\textbf{Component}& \textbf{Amount}\\                           
\endfirsthead
\rowcolor{white}
\caption{Continuation}
\\
\hline
\textbf{Component}& \textbf{Amount}\\ 
\endhead
\hline
\endfoot
GTE-Buffer & 25mM Tris pH 8.0,  10mM EDTA pH 8.0, 50mM Glucose, 24 µg/ml RNAse A \\
KaC-Buffer &  \\
Lysis-Buffer (50ml) &  1\% SDS, 0.2M NaOH      \\
\end{longtable}

\subsection{Medium scale Plasmid DNA preparation}

When larger quantity of DNA product from recombinant E. coli bacteria was needed, growth and extraction in 200 ml format was used. 200 ml of LB medium were supplemented with the appropriate antibiotics to select for the desired recombinant strains and inoculated with a single clonal bacterial culture. The fluid culture was incubated overnight at 37\C\ while rotating for a maximum of 16 hours. It was then centrifuged for 20 minutes at 5000 x g and 4 \C. Preparation was done with the Genopure Midi prep Kit by Roche and with the manufacturers protocol recommendations. The supernatant was removed and the bacterial pellet suspended in 8ml suspension buffer with added RNAse. 8 ml lysis buffer were added and the tube inverted 6-8 times, then incubated 5 minutes at room temperature. 8 ml cold neutralisation buffer were added and incubation continued for another 15 minutes on ice. Afterwards, the lysate was centrifuged for 30 minutes at 18,000 x g at 4\C. Supplied eluation collumns were treated with 2.5 ml equilibration buffer and the lysate supernatant was given into the column. Flow-through was discarded and the column matrix treated with 4 ml washing buffer 2 times. The column was placed on a fresh falcon and the DNA eluted from the matrix with 5 ml eluation buffer. DNA was precipitated with 3.6 ml isopropanol. The additional washing and concentration steps are identical with protocol \ref{freezy-squeezy}.

\subsection{Transformation of DH10\(\beta\) cells by electroporation}

Introduction of foreign DNA into a recombinant E. Coli strain was done by electroporation. Cells were thawed, cuvettes and plasmid DNA also chilled on ice and cuvettes were thoroughly dried before use. Immediately after thawing, 5 ng of Plasmid DNA was added to the cell suspension, the mixture transferred to the cuvette and electroporation was done in  at 1,350 V, 10 \(\mu\)F and 600 Ohms. Optimal electroporation times was aimed to be 5 ms. After the device was finished, the cells were mixed with 1 ml warmed LB medium (37 \C) without antibiotics and shaken for approximately one hour before streaking on agar dishes.

\section{Molecular biological Methods}

\subsection{DNA storage}

Oligonucleotides were stored at -20\(^\circ\mathrm{C}\) in \ce{H2O}. Plasmids and bacmids were stored in TE Buffer with a pH of 8.0 at  4\(^\circ\mathrm{C}\). 

\begin{longtable}{C{0.93\linewidth}}
\rowcolor{white}
\caption{TE Buffer composition}
\label{TE buffer composition} 
\\
\hline
\textbf{Reagents and Amounts} \\
\endfirsthead
\endfoot
\mbox{10 mM Tris-HCL}, \mbox{1 mM EDTA},  pH 8.0 \\
\hline
\end{longtable}


\subsection{Polymerase-Chain-Reaction}

PCR is a classic method to produce large amounts of existing genomic material or detect a suspected nucleic acid. For the reaction, DNA primers are needed that originate in sequence from the flanks of the nucleic acids to be amplified. The amplification reaction is catalyzed by a heat-resistant polymerase enzyme. Chronologically, PCR starts with the denaturation of the DNA, then proceeds with the annealing of the primers to the template, and finally with the elongation. For diagnostic use, where the interest was in the size of the resulting PCR products, the GoTaq polymerase was used in combination with the green GoTaq buffer. When preparing material for sequencing, high-fidelity Q5-polymerase with the supplied buffer was used. According to the manufacturers' recommendations, the different polymerases required slightly different reaction conditions, which were applied accordingly. For each reaction, a 50 \ul PCR tube was used.  The annealing time was adjusted to the specific primers used in each reaction. Tables \ref{tab:PCR scheme for GoTaq Polymerase} to \ref{Basic pipetting scheme for Q5 polymerase} show the programs and pipetting instructions for each polymerase. Table \ref{tab:primers} shows the primers used in this project.


\begin{table}[H]
    \centering
\caption{PCR scheme for GoTaq Polymerase}
\label{tab:PCR scheme for GoTaq Polymerase}
    \begin{tabular}{>{\raggedright\arraybackslash}p{0.21\linewidth}>{\centering\arraybackslash}p{0.21\linewidth}>{\centering\arraybackslash}p{0.21\linewidth}>{\raggedright\arraybackslash}p{0.21\linewidth}}
    \hline
      \textbf{Step}   & \textbf{Temperature}  & \textbf{Time (min)}  & \textbf{Steps}\\
    Denaturation & 98 & 2 & \\
 \hline 
    Denaturation &  98\C& 0.5 & \multirow{3}{=}{\cellcolor{white}x35}\\
    Annealing  & X & 0.75 \\
    Elongation & 72\C & 0.5/kb  \\
\hline
    Elongation & 72\C & 15 &  \\
         \hline
      \end{tabular}
\end{table}

\begin{table}[H]
    \centering
\caption{PCR scheme for Q5 Polymerase}
\label{tab:PCR scheme for Q5 Polymerase}
    \begin{tabular}{>{\raggedright\arraybackslash}p{0.21\linewidth}>{\centering\arraybackslash}p{0.21\linewidth}>{\centering\arraybackslash}p{0.21\linewidth}>{\raggedright\arraybackslash}p{0.21\linewidth}}
    \hline
      \textbf{Step}   & \textbf{Temperature}  & \textbf{Time (min)}  & \textbf{Steps}\\
    Denaturation & 98 & 2 & \\
 \hline 
    Denaturation &  98\C& 0.5 & \multirow{3}{=}{\cellcolor{white}x35}\\
    Annealing  & X & 0.75 \\
    Elongation & 72\C & 1/kb  \\
\hline
    Elongation & 72\C & 20 &  \\
         \hline
      \end{tabular}
\end{table}

\begin{longtable}[H]{L{0.45\linewidth} R{0.45\linewidth}}
\rowcolor{white}
\caption{basic pipetting scheme for GoTaq polymerase}
\label{Basic pipetting scheme for GoTaq polymerase}
\\
\hline
\textbf{Component}& \textbf{Volume}\\                           
\endfirsthead
\rowcolor{white}
\caption{Continuation}
\\
\hline
\textbf{Component}& \textbf{Volume}\\ 
\endhead
\hline
\endfoot
x5 green GoTaq Buffer & 10 \ul \\
dNTPs & 1 \ul \\
Forward primer (10 \(\mu M\)) & 1 \ul \\
Reverse primer (10 \(\mu M\)) & 1 \ul \\
Template DNA & 50 ng \\
GoTaq polymerase & 0.25 \ul \\
\ce{H2O} & 36,25 \ul \\
\end{longtable}

\begin{longtable}[H]{L{0.45\linewidth} R{0.45\linewidth}}
\rowcolor{white}
\caption{basic pipetting scheme for Q5 polymerase}
\label{Basic pipetting scheme for Q5 polymerase}
\\
\hline
\textbf{Component}& \textbf{Volume}\\                           
\endfirsthead
\rowcolor{white}
\caption{Continuation}
\\
\hline
\textbf{Component}& \textbf{Volume}\\ 
\endhead
\hline
\endfoot
x5 Q5 Buffer & 10 \ul \\
dNTPs & 1 \ul \\
Forward primer (10 \(\mu M\)) & 2.5 \ul \\
Reverse primer (10 \(\mu M\)) & 2.5 \ul \\
Template DNA & 300 ng \\
Q5 polymerase & 0.5 \ul \\
\ce{H2O} & 33 \ul \\
\end{longtable}

\begin{longtable}[H]{L{0.45\linewidth} R{0.45\linewidth}}
\rowcolor{white}
\caption{Primers used in this project}
\label{tab:primers}
\\
\hline
\textbf{Name}& \textbf{Sequence}\\                        
\endfirsthead
\rowcolor{white}
\caption{Continuation}
\\
\hline
\textbf{Name}& \textbf{Sequence}\\ 
\endhead
\hline
\endfoot
Left junction fwd 1 (711) &  \footnotesize \uppercase{gaaccaatatgaactgcaggc} \\
Left junction fwd 2 (713) & \footnotesize \uppercase{ctaaaagagcgcctacatgg}  \\
Right junction rev 1 (716) & \footnotesize \uppercase{gaggattctgttcagagacctg}   \\  
Right junction rev 2 (718) & \footnotesize \uppercase{AAATGGTGAGAGCCATCCAC}  \\
E1a fwd (763) & \footnotesize \uppercase{tacccgccgtcctaaaatgg}   \\  
E1b rev (766) & \footnotesize \uppercase{ggaacagcgggtcagtatgt}  \\          
&   \\  
&   \\
\end{longtable}

\newpage

\subsection{Agarose Gel Electrophoresis}
\label{electrophoresis}
DNA Gel electrophoresis utilises the proportionally increasing, length-dependent travel speed of negatively charged DNA molecules through a gel matrix with applied electric current. Travel over time can be influenced by varying the pore size of the gel matrix or the voltage applied to the gel. In this work, 0.8\% agarose gels were used. The gels were prepared with the addition of one drop ethidium bromide, a fluorescent DNA intercalating agent, for identifying fragment sizes. The runtime varied from 30 to 50 minutes at a voltage of  5 V/cm, varying slightly depsending on expected fragment sizes. A Gel Doc EZ camera system was used to document the finished gels. 

\subsection{Elution of DNA from agarose gels}
\label{freezy-squeezy}
When further use of DNA fragments after electrophoresis (\ref{electrophoresis}) was intended, the position of the desired band was confirmed by inspection under UV-light and the gel was cut out with a scalpel and minced finely. a 1:1 ratio by weight of phenol was added, the reaction tube shaken vigorously for one minute and frozen in liquid nitrogen for an additional minute before being thawed completely in hand. The tube was centrifuged for 15 minutes at 19,000 x g (same g s for all centrifuging steps in this protocol). the upper aqueous phase was separated into a new reaction tube. a 1:1 ratio of chloroform-isoamylalcohol 24:1 was added, shaking and spinning were repeated like above. the upper phase was again separated and 0.1 volumes 3M NaAc pH5.2 along with 0.7 volumes isopropanol 100\% were added. This was left in the freezer at -20\C\ for 30 minutes or alternatively overnight. Centrifuging was repeated and the resulting pellet solved in 150 \ul TE buffer. To further enhance DNA purity, precipitation was repeated with 0.1 volumes NaAc and 2.5 volumes Ethanol 100\%. The Centrifuged pellet was then washed with Ethanol 70\%, centrifuged and solved in 30 \ul TE buffer.

\subsection{px330 expression plasmid and BbsI restriction enzyme}

The plasmid px330, an aliquot of which was kindly given to us by Dr. Eric Schulz from UWH, carries the Cas9-sequence, barring the sequence of the guide-rna, and a puromycin resistance encoding sequence. \\
BbsI is a restriction  enyzme originating from \textit{bacillus brevis}, recognising a single six-nucleotide motif. The px330 plasmid has two matching bbsI-sites, that are cut away after bbsI-restriction and insertion of a gRNA-sequence.

\subsection{DNA oligo preparation}
Single-stranded DNA oligos with the gRNA-sequences were ordered in complimentary forward and reverse orientation from ThermoFisher. The stocks were diluted to a concentreation of 1 \(\mu\)g/\(\mu\)l. 1 \(\mu\)g each of  matching forward and reverse oligos were pipetted ad 50 \ul TE 8.0 in a reaction tube and boiled in  water on a heating plate for 15 minutes. the heating plate was then turned off and the oligos were left to cool overnight at room temperature. Table \ref{oligo seq} shows the sequences of the gRNA oligos

\begin{longtable}{L{0.45\linewidth} R{0.45\linewidth}}
\rowcolor{white}
\caption{gRNA oligo sequences}
\label{oligo seq}
\\
\hline
\textbf{Name}& \textbf{Sequence}\\                           
\endfirsthead
\rowcolor{white}
\caption{Continuation}
\\
\hline
\textbf{Name}& \textbf{Sequence}\\ 
\endhead
\hline
\endfoot
gRNA 1 fwd &  \footnotesize{CACCGGGCGGAACACATGTAAGCGA} \\
gRNA 1 rev &  \footnotesize AAACTCGCTTACATGTGTTCCGCCC \\
gRNA 2 fwd &  \footnotesize CACCGGTAAACGGTCAAAGTCCCCG \\
gRNA 2 rev & \footnotesize  AACCGGGGACTTTGACCGTTTACC  \\
gRNA 3 fwd & \footnotesize CACCCCGTCGGTATTATCGCCGGG  \\
gRNA 3 rev &  \footnotesize AAACCCCGGCGATAATACCGACGGC \\
gRNA 4 fwd & \footnotesize CACCTGTACAACATTCATTCCCGA  \\
gRNA 4 rev & \footnotesize AAACCGTTCAGGTTCACCTTGGACC \\
gRNA 5 fwd & \footnotesize CACCGTCCAAGGTGAACCTGAACG  \\
gRNA 5 rev & \footnotesize AAACCGTTCAGGTTCACCTTGGACC \\
gRNA 6 fwd & \footnotesize CACCGTGTTTACCGCCACACTCGCA \\
gRNA 6 rev & \footnotesize AAACTGCGAGTGTGGCGGTAAACAC \\
\end{longtable}

\subsection{Restriction-ligation cloning of DNA oligos into plasmid backbones}

To linearize the 10 ng of px330 plasmid was incubated with 20 units BbsI ad 100 \ul \ce{H2O} supplemented with appropriate measure of NEB 2.1 enzyme buffer at 37 \C for one hour. Then, the whole reaction solution was subjected to gel electrophoresis like described under \ref{electrophoresis} to seperate successfully linearized plasmid and purified from the finised gel run like shown in \ref{freezy-squeezy}. The plasmid DNA concentration of the eluate was then measured and amounts for the ligation mixed like in table \ref{ligation reaction pipetting scheme}. Table  

\begin{longtable}[ht]{L{0.45\linewidth} R{0.45\linewidth}}
\rowcolor{white}
\caption{ligation reaction pipetting scheme}
\label{ligation reaction pipetting scheme}
\\
\hline
\textbf{Component}& \textbf{Amount}\\                           
\endfirsthead
\rowcolor{white}
\caption{Continuation}
\\
\hline
\textbf{Component}& \textbf{Amount}\\ 
\endhead
\hline
\endfoot
10x T4 ligase Buffer & 1 \ul \\
lin. px330 plasmid & 100 ng \\
ds gRNA oligo & 0.55 \ul \\
T4 ligase & 0.5 \ul \\
\ce{H2O} & 5.95 \ul \\
\end{longtable}

First, only \ce{H2O} and the plasmid were mixed, incubated for 10 minutes at 42\C \ to linearize any supercoiled DNA structures that would impede enzyme efficacy. The mixture was then cooled for 5 minutes on ice. Buffer, ligase and oligos were added to the reaction tubes and the ligation was catalyzed at 16\C \ overnight in a thermocycler. 

\newpage

\section{Cell biological Methods}

\subsection{Eucaryotic Cell culture}
\label{Cell culture}

A549 and HEK293 cells were used. They were cultured in \mbox{15 cm dishes} and kept in MEM medium supplemented with 10\% FBS and 1\% penicillin / streptomycin. HEK cells were split at around 80-90 percent confluency. To split the cells, they were washed with PBS, treated with Trypsin for 1 minute before stopping the reaction with unsupplemented MEM medium and then repeatedly pipetting up and down. Depending on the confluency, the splitting ratio was 1:8 to 1:10. A549 cells were treated in the same manner, except cells were trypsinised for 3 minutes and split at a ratio of 1:6 to 1:8. All reagents used were warmed to 37\ \(^\circ\mathrm{C}\) before being used. Dishes were stored at 37\ \(^\circ\mathrm{C}\) in an incubator with an atmosphere of 5\% \ce{CO2}. 


\subsection{Cell Counting}
\label{cell counting}

To determine cell numbers for usages where exact quantities are needed, cells were removed from their dishes and seperated like described above and resuspended in MEM medium. An improved Neubauer chamber was loaded with the cell suspension and all four squares were counted under the light-microscope. the mean of those squares was then taken to calculate cell count like \[[Cells/ml] = [Mean/squares]  1 * 10^4\] This value was then used to calculate the desired volume to use in cell seeding.

\subsection{Transfection of HEK293 cells with Cas9 Plasmids}

24 hours prior to transfection HEK293 cells were seeded in 6cm cell culture dishes counted to an amount of 5x10\(^5\) per dish. Dishes were filled with 2 ml of medium. All gRNA constructs were transfected in parallel. Per gRNA, 500 ng were used, 3 \(\mu\)g for each cell culture dish in total. Transfection was prepared by mixing the correct amounts of plasmid solution and filling up woth 150 mM NaCl to a volume of 250 \ul per transfection aliquot. The appropriate measure of Polyethylenimine (PEI) was calculated with the following formula: \[PEI \ [\mu l] = (\mu g \ DNA) \ x \ 12\]
As calculated with the used DNA amounts, 36 \ul PEI were diluted to a volume of 250 \ul 150 mM NaCl for each transfection. PEI is a molecule that forms complexes with DNA of positive net charge, allowing for higher rates of intracellular DNA delivery that would normally be hindered by the opposing negative charges of the DNA and the cellular membrane. The NaCl solution containing the PEI was given into the DNA aliquots carefully, aditionally pipetted up and down three times and then incubated at room temperature for 10 minutes. Finally the mixture was added to the cell dishes after they were washed with PBS and supplemented with fresh Medium. After at least 2 hours in the incubator, medium was filled up to a volume of 6 ml per dish.

\subsection{Extraction of genomic DNA}

To yield high purity genomic DNA, to use for sequencing or PCR-assays, the phenol-chloroform extraction method was used in this project. Cells were removed from 15 cm culture dishes as described in \ref{Cell culture}. for one 1.5 ml microcetrifuge tube, half of the cells from a confluent 15 cm dish were used. Medium was removed by centrifuging at 300 x g for 4 minutes and the cell pellet was re-suspended in 200 \ul PBS. 200 \ul Lysis-Buffer was added and the suspension was mixed thoroughly before adding 30 \ul SDS 10\% and 20 \ul proteinase K. The mixture was shaken overnight at 55\C and 1,200 rpm.
After overnight shaking, approximately 16 hours later, 20 \ul RNAse A was added and the solution was incubated for a further 30 minutes at 37\C/900rpm. 470 \ul Phenol was added and the tube shaken by hand for one minute, then centrifuged at 11,000 x g for 10 minutes and resulting supernatant was carefully transferred into another microcentrifuge tube and 450 \ul Phenol was added. Shaking, centrifuging  and transferring were repeated twice as above, swapping phenol for 430 \ul chloroform/isoamylalcohol in a ratio of 24:1 before the second repetition. 40 \ul \ce{NaAc} and 1ml Ethanol 100\% were added and the tube was inverted until precipitated DNA could be visually confirmed in the tube. centrifuging at full speed for 20 minutes at room temperature yielded a DNA pellet which was separated from the supernatant and washed with 500 \ul Ethanol 50\%, centrifuged at full speed for 5 minutes, air-dryed for one minute and suspended in 50 \ul TE-Buffer before being stored at 4\C \ overnight to re-enter solution. DNA quantity in the solution was then measured like described in \ref{OD measure}.  

\begin{longtable}{C{0.93\linewidth}}
\rowcolor{white}
\caption{Lysis buffer composition}
\label{Lysis buffer composition} 
\\
\hline
\textbf{Reagents and Amounts} \\
\endfirsthead
\hline
\endfoot
\mbox{10 mM Tris ultrapure}, \mbox{10 mM EDTA},  \mbox{0,5 \% SDS}
\end{longtable}
\subsection{genomic DNA quantification}
\label{OD measure}
To approximate the DNA concentration in a given solution, OD$_{\textnormal{\scriptsize 260}}$ was measured using a photometer, which was calibrated using a blank solution. DNA concentration was then calculated by the device and displayed as ng/\(\mu\)l.

\begin{longtable}{L{0.45\linewidth} R{0.45\linewidth}}
\rowcolor{white}
\caption{Composition of solutions for DNA quantification}
\label{Composition of solutions for DNA quantification}
\\
\hline
\textbf{Solution}& \textbf{Composition}\\                           
\endfirsthead
\rowcolor{white}
\caption{Continuation}
\\
\hline
\textbf{Solution}& \textbf{Composition}\\ 
\endhead
 blank solution& 1 \ul TE 8.0, 69 \ul \ce{H2O} \\
 measured solution& 1 \ul DNA elution, 69 \ul \ce{H2O} \\
\hline
\endfoot


\end{longtable}

\subsection{Whole-cell protein extraction}

To prepare protein extract for SDS-page, cells were seeded in 12 well cell culture dishes and grown to $\simeq$60\% confluency. Cells were then washed with PBS and treated with 300 \ul lysis buffer for 10 minutes at room temperature. 100 \ul TCA 100\% was then added and mixing was done by gently shaking the culture dish. The cloudy supernatant was transferred to a 1.5 ml reaction tube, leaving a small white sphere of denatured DNA in the dish well. the tube was centrifuged for 5 minutes at 5,000 x g, room temperature and the pellet washed with TCA 2.5\%. Centrifuging was repeated like above and the pellet washed with 20 \ul tris-base and vortexed briefly. The tube was incubated for 30 minutes at room temperature and then suspended in 20 \ul \ce{H2O} to re-enter solution. Protein extract was stored at -20\C \ in the freezer.


\subsection{SDS-Page}
\label{SDSpage}

\begin{longtable}{L{0.45\linewidth} R{0.45\linewidth}}
\rowcolor{white}
\caption{SDS-Page Buffers and Solutions}
\label{SDSBuffers}
\\
\hline
\textbf{Buffer}& \textbf{Composition}\\                           
\endfirsthead
\rowcolor{white}
\caption{Continuation}
\\
\hline
\textbf{---}& \textbf{---}\\ 
\endhead
\hline
\endfoot
10\% Separating Gel & 6 ml \ce{dH2O}, 3.75 ml 4x Separating Gel Buffer pH 8.8, 5 ml 30\% PAA (29:1), Combine and use 5 ml of the solution, 50 µl 10\% APS, 5 µl TEMED \\
5\% Stacking Gel &  3.4 ml dH2O, 1.5 ml 4x Stacking Gel Buffer pH 6.8. 1 ml 30\% PAA (29:1), Combine and use 2.5 ml of the solution, 30 µl 10\% APS, 3 µl TEMED  \\
Seperating Gel Buffer & \\
Stacking Gel Buffer & \\ 
\end{longtable}
The separation gel was first poured into the a gel preparation chamber. Before pouring the stacking gel, a few milliliters of isopropanol were added to remove any bubbles. after 10 minutes, the isopropanol was removed and the stacking gel poured. Protein samples were mixed with SDS loading buffer and denatured for 10 minutes at 75\C. Voltage was set to 4V/cm and increased to 5 V/cm once the dye front reached the Separating gel. The gel-chamber was placed on ice for the whole run to protect the gel matrix integrity. 

\subsection{Western Blot}
\label{Western blot}

\begin{longtable}{L{0.45\linewidth} R{0.45\linewidth}}
\rowcolor{white}
\caption{Western Blot Buffer composition}
\label{WBbuffercomp}
\\
\hline
\textbf{Buffer}& \textbf{Composition}\\                           
\endfirsthead
\rowcolor{white}
\caption{Continuation}
\\
\hline
\textbf{Buffer}& \textbf{Composition}\\ 
\endhead
\hline
\endfoot
10x Tris-Glycin buffer & 30 g/l Tris pH 8,3, 144 g/l Glycin \\
1x Transfer buffer (1L) & 100 ml 10x Tris-Glycin, 200 ml Methanol, 700 ml \ce{H2O} \\
10x TBS	(1L) & 100 mM Tris, 1.5 M NaCl, 60 ml HCl ad 1L \ce{ddH2O} ph 7.4 \\
1x TBS-T & 450 ml \ce{H2O}, 50 ml 10x TBS, 250 µl Tween-20 \\
Blocking buffer & 5\% BSA in TBS-T\\
Primary antibody & 1:1,000 in blocking buffer\\
Secondary antibody & 1:15,0000 in blocking buffer\\
\end{longtable}

Gels resulting from SDS-PAGE (\ref{SDSpage}) were subjected to lateral electric current of  5 V/cm for one hour in an ice-cooled blotting chamber. The gel was placed adjacent to a nitrocellulose membrane, thus allowing travel of the proteins along the current onto the secondary membrane. After transfer, the membrane was washed in TBS, shortly air-dryed and placed in blocking buffer overnight at 4\C. After blocking, the membrane was incubated with the primary antibody for 1 hour at room temperature and then rinsed 5 times in TBS-T. After rinsing, incubation with the secondary antibody was performed for 45 minutes at room temperature and rinsing was repeated as done above. The membrane was then briefly dried and signals were detected in the Odyssey CLx Infrared Scanner.


\subsection{Nuclear staining of E1a}
\label{E1 nuclear stain}

\begin{longtable}{L{0.45\linewidth} R{0.45\linewidth}}
\rowcolor{white}
\caption{Nuclear stain Buffer composition}
\label{Nuclear stain buffers}
\\
\hline
\textbf{Buffer}& \textbf{Contents}\\                           
\endfirsthead
\rowcolor{white}
\caption{Continuation}
\\
\hline
\textbf{Buffer}& \textbf{Contents}\\ 
\endhead
\hline
\endfoot
Fixation buffer & 4\% PFA in PBS \\
Permeabilization buffer & \mbox{0.1\% Triton X-100}, 0.68 \(\mu\)M EDTA, 1\% BSA in PBS\\
Incubation buffer & 0.5\% BSA in PBS \\

\end{longtable}


Cells were prepared for seeding like described in \ref{cell counting}. 5x10\(^5\) cells were used per well of a conical 96-well plate. after detatching, cells were centrifuged and re-suspended in 100 \ul fixation buffer at room temperature for 15 minutes, then centrifuged again. All centrifuging steps, unless otherwise stated, are performed at 4\C, 300 x g, for 5 Minutes. the cells were mixed with 50 \ul permeabilization buffer and incubated on ice, 10 minutes, then centrifuged, washed in 100 \ul incubation buffer, centrifuged and 50 \ul incubation buffer with the diluted first antibody was added. The incubation time for the first antibody was 45 minutes at 4 \C, re-suspending the cells once to ensure  even distribution in the solution. After incubation, cells were spun, immediately followed by mixing with 50 \ul incubation buffer with the diluted second antibody. incubation times and temperature was repeated like before and finally, cells were washed with incubation buffer and re-suspended in 150 \ul of incubation buffer and transferred to a 96 well FACS plate and proceeding with the measurements.


\subsection{Ad-vector replication Assay}
\label{production assay}

To indirectly assess E1-expression, E1-knockdown HEK293 clones were infected with a recombinant Ad5 vector expressing eGFP under control of the MLP promoter, which is active only when replication occurs. Since replication of E1-deficient Ad5 vectors is dependent on cellular E1 transcomplementation, changes in E1 expression can be extrapolated by measuring eGFP flourescence with flow cytomerty. Cells were counted to 2x10\(^4\)/100 \ul medium and seeded in 96-well cell culture dishes. 24 hours later, cells were infected with a MOI of 300 in a volume of 20 \ul cell culture medium. 28 hours post infection. cells were fixated and prepared and measured like described under \ref{FC}

\subsection{Flow Cytometry with Eucaryotic Cells}
\label{FC}
In this project, methods \ref{E1 nuclear stain} and \ref{production assay} used flow cytometry (FC) measurements to generate Data. Prior to FC, cells had the medium removed, were removed from dish binding with trypsin, after 1 minute the reaction was halted with medium. The cells were spun for 5 minutes at 400 x g, room temperature, the pellet re-suspended in PFA 4\% and incubated for 15 minutes. Centrifuging was repeated and the cells suspended in 130 \ul PBS and transferred to a FC plate. In FC settings, gating for single cells was done for all measurements and data was extracted as .csv.

\subsection{Statistical analysis}

\subsection{DNA Sequencing}

DNA sequencing was done using sanger-sequencing by the Microsynth AG, DNA was prepared as per the requirements of the service provider: For plasmid DNA teh preparation was 50 ng/\ul \ in a total Volume of 12 \ul \ \ce{H2O}. For PCR products, varying concentrations depending on the fragment length were used.



